\documentclass[10pt,a4paper]{amsart}
\usepackage[margin=19mm]{geometry}
\usepackage{amsmath,amssymb,mathtools}
\usepackage[scr=rsfs,cal=euler]{mathalfa}
\usepackage{xfrac}
\usepackage[unicode,hidelinks,bookmarks=false]{hyperref}
\newcommand{\ie}{i.e.\ }
\newcommand{\cf}{cf.\ }
\newcommand{\resp}{respectively\ }
\newcommand{\Subsection}{\subsection}
\providecommand{\nobreakdash}{\nobreak\hbox{-}}
\setlength{\emergencystretch}{2em}
\multlinegap0pt
\usepackage{amssymb}
\usepackage{enumerate}
\usepackage[normalem]{ulem}
\newtheorem{proposition}{Proposition}
\newtheorem{theorem}{Theorem}
\newcommand{\C}{\mathbb C}
\newcommand{\maC}{\mathcal C}
\newcommand{\maD}{\mathcal D}
\newcommand{\maH}{\mathcal H}
\newcommand{\maK}{\mathcal K}
\newcommand{\maL}{\mathcal L}
\newcommand{\maP}{\mathcal P}
\newcommand{\oid}{\operatorname{Id}}
\title{Fredholm Criteria for $G$-pseudodifferential Operators}
\author{Alexandre Baldare, Anton Yu. Savin, Elmar Schrohe}
\date{Source: arXiv:2605.14778v1 (CC BY 4.0)}
\begin{document}
\maketitle

\noindent\emph{Source and licence.} Fredholm Criteria for $G$-pseudodifferential Operators. Version arXiv:2605.14778v1 (CC BY 4.0), distributed by arXiv under the Creative Commons Attribution 4.0 International licence, http://creativecommons.org/licenses/by/4.0/, as recorded in the arXiv licence metadata for this exact version and shown on the abstract page https://arxiv.org/abs/2605.14778v1. Full licence notice: \texttt{licenses/arxiv-2605.14778-CC-BY-4.0.txt}.\par\smallskip

\noindent\textit{Editorial scope.} Counted unit: the article's theorem thm.Simonenko.Gop (source0.tex lines 576--584). The article's complete original proof of it is delivered with it (source0.tex lines 586--660, one \texttt{proof} environment of 75 lines). No mathematics is written, repaired or re-derived: the statement and the proof are the article's own lines, in the article's own notation, and no retained line is substituted or re-flowed.  Retained uncounted as the source-local context the proof needs: lines 285--312; lines 499--508.  Nothing was omitted from any retained span, so the package carries no omission record.  No retained line contains a citation, so the wrapper carries no bibliography and no bibliography entry.  No retained line contains a figure environment, an image-inclusion command or drawing code, so no image is delivered and the package declares no asset dependency.  Editorial formatting only, in addition: an AMS \texttt{amsart} preamble that restates the article's own declarations and macros used by the retained lines, namely the environments \texttt{proposition}, \texttt{theorem} on separate counters, together with the standard packages those lines need.  The retained environments print in this document as Proposition 1, Theorem 1 in order of appearance, whereas the article prints the same items with the numbering of its own full document.  The complete native source of the article as distributed in the arXiv e-print for this exact version is bundled unchanged as \texttt{sources/200612/source0.tex}.
\begin{equation*}
\sigma_0(\hat{\maD})(\xi) : E_\xi \otimes \C[\Gamma] \rightarrow E_\xi \otimes \C[\Gamma]  
\end{equation*}
is sufficient for the Fredholm property of $\maD$ and that it is also necessary,  if the action of $\Gamma$ is efficient. If the action is efficient, a minimal isotropy group is trivial, and we recover Antonevich's result. On the other hand, if a minimal isotropy group is not trivial, we  obtain a new necessary and sufficient condition for the Fredholm property.


\medskip

\noindent
{\em{Acknowledgements.}} The authors would like to thank RFBR,  project number 21-51-12006, and DFG, project SCHR 319/10-1 as well as the DFG Priority Programme 2026 `Geometry at Infinity'   
    for their support.
 The first named author would like to express his gratitude to V. Nistor for useful discussions, suggestions and encouragement during the redaction of this paper.   
 He also wishes to dedicate this work to the memory of M. Benameur, his doctoral advisor, who passed away in September 2025 and who taught him to reach for the stars -- and $C^*$-algebras were the closest ones.  
  


\section{Preliminaries}\label{sec.preliminaries}

Let $G$ be a compact Lie group which acts smoothly on a compact Riemannian manifold $M$ without boundary. 
Without loss of generality, we assume that $M$ is endowed with a $G$-invariant metric and that $G$ acts by isometries.
Let $E^\pm \rightarrow M$ be $G$-equivariant vector bundles equipped with $G$-invariant Hermitian metrics.
We shall denote by $\psi^m(M,E^+,E^-)$ the space of  
classical pseudodifferential operators of order $m$, acting between
sections of the $G$-equivariant vector bundles $E^\pm \rightarrow M$.
As pointed out above, we are interested in operators of the form
\begin{equation*}
P + \maD : C^\infty(M,E^+) \rightarrow C^\infty(M,E^-),\qquad (P + \maD) s(x)=Ps(x) + \int_G D_g (T_gs) (x) dg,
\end{equation*}

\begin{proposition}[General Simonenko localization principle]\label{prop.gen.Simonenko} 
%AS \  \begin{enumerate}\renewcommand\labelenumi{\rm (\alph{enumi})}
% \item \item
Let $P \in \Psi_T(\maH)$.
%Assume that $M_\theta$ goes strongly to zero when the support of $\theta$ becomes small. 
Assume that $\maC(T)\hookrightarrow \maL(\maH)$ has the property of strong convergence to $0$.
Then
$P$ is locally invertible on $T$ if, and only if, $P$ is Fredholm.
%\end{enumerate}
\end{proposition}

\begin{theorem}[Simonenko's principle for $G$-operators] \label{thm.Simonenko.Gop}
Let $\maD \in \Psi_{G}(M,E)$ be a $G$-operator as in Section \ref{sec.preliminaries}.
%\begin{enumerate}
%\item Then $M_{\chi}(\oid + \maD)M_{\chi} \in \Psi_{M\smallsetminus \maP/G}(L^2(M \smallsetminus \maP,E))$.
%\item Furthermore, $M_{1-\chi} \maD M_{1-\chi} \in \maK(L^2(\maP,E))$.
%\item We have that $\oid + \maD$ is locally invertible if, and only if, it is Fredholm.
%\end{enumerate}
Then $\oid +\maD$ is Fredholm if, and only if, $p \maD p + \oid_{L^2(M\smallsetminus \maP , E\vert_{M\smallsetminus \maP})}$ is locally invertible.
\end{theorem}

\begin{proof}
%1. Let $\maD \in \Psi_{G}(M,E))$ and $\phi,\psi \in C(M)^G$ be $G$-invariant functions with disjoint supports. 
%Notice that $\forall g \in G$, 
%we have $T_g M_\psi=M_\psi T_g$. Therefore, it follows
%that 
%\begin{equation*}
%M_\phi (\oid + \maD)M_\psi = M_\phi \maD M_\psi 
%%= \int_G M_\phi D_g T_g M_\psi  \ dg
%= \int_G M_\phi D_g  M_\psi  T_g \ dg \in \maK(L^2(M,E)), 
%\end{equation*}
%because $M_{f_1} P M_{f_2}$ is compact $\forall P \in \Psi(M,E)$ and $f_1,f_2 \in C(M)$ with disjoint support.
%
%2. This follows directly from item 1 and Proposition \ref{prop.gen.Simonenko}.

Recall that $p=M_{\chi_{M\smallsetminus \maP}}$, $\oid-p = M_{1-\chi_{M\smallsetminus \maP}}$ and 
let us write 
\begin{equation*}
\maD= p \maD p + (\oid-p) \maD (\oid-p) +p \maD (\oid-p) + (\oid - p) \maD p.
\end{equation*}
The operators $(\oid-p) \maD (\oid-p)$, $p \maD (\oid-p)$ and $(\oid - p) \maD p$ are compact.
Indeed, let 
$\phi,\psi \in \maC(M)^G$ be $G$-invariant functions with disjoint supports. 
Then  $T_g M_\psi=M_\psi T_g$ for all $g\in G$ and therefore
\begin{equation*}
M_\phi \maD M_\psi 
= \int_G M_\phi D_g  M_\psi  T_g \ dg \in \maK(L^2(M,E)), 
\end{equation*}
because $M_{f_1} P M_{f_2}$ is compact whenever $ P \in \Psi(M,E)$ and $f_1,f_2 \in \maC(M)$ have disjoint supports.
This implies that $p \maD (\oid-p)$, and $ (\oid-p) \maD p$ are compact operators.

To show that $(\oid-p)\maD (\oid-p) $ is a compact operator, it is enough to show that 
$(\oid-p) \maD (\oid-p)\in \maK(L^2(\maP,E\vert_\maP))$. 
Recall that $\maP=\sqcup C_i$ is a finite union of clopen orbits $C_i \cong G/G_{x_i}$ with suitable $x_i\in M$ 
and that $E\vert_{C_i} \cong G \times_{G_{x_i}} E_{i}$,
where $E_{i}$ is a unitary representation of $G_{x_i}$.    
Let $\phi_i$ be the characteristic function of $C_i$. Then we can write
\begin{equation*}
M_{1-\chi} \maD M_{1-\chi} = \sum_i M_{\phi_i} \maD M_{\phi_i} + \sum_{i\neq j} M_{\phi_i} \maD M_{\phi_j}.
\end{equation*} 
From the previous discussion we know that $M_{\phi_i} \maD M_{\phi_j}$ is compact if $i \neq j$ 
and therefore $\sum_{i\neq j} M_{\phi_i} \maD M_{\phi_j}$ is compact.
Let us show that also $M_{\phi_i} \maD M_{\phi_i} $ is compact.
Notice that $M_{\phi_i} \maD M_{\phi_i} = \int_G M_{\phi_i} D_g M_{\phi_i} T_g \ dg$ is a $G$-pseudodifferential operator
on $C_i \cong G/G_{x_i}$. By definition, $M_{\phi_i} \maD M_{\phi_i}$ is the limit of sums of operators of the form 
$B \circ T_\varphi$, where $B$  is a bounded operator on $L^2(C_i,E\vert_{C_i})$, 
$T_\varphi:=\int_G \varphi(g) T_g \ dg$ and $\varphi \in C^\infty(G)$.
Since %Recall that $C_i\cong G/G_i$ and that 
$E\vert_{C_i} \cong G \times_{G_{x_i}} E_{i}$ %, where $E_i$ is a unitary representation of $G_i$. Therefore, if 
and $\varphi \in C^\infty(G)$  the operator 
$ \hat{T}_\varphi:=\int_G \varphi(g) T_g \ dg : L^2(G,E_{i}) \rightarrow L^2(G,E_{i})$ is compact, 
because its Schwartz kernel is given by the smooth function $k_\varphi(g,h):=\varphi(gh^{-1})$. 
It follows that the restriction $\hat{T}_\varphi^{G_{i}}$ of $\hat{T}_\varphi$ to 
$L^2(G,E_{i})^{G_{x_i}}\cong L^2(G/G_{x_i},G \times_{G_{i}} E_{i})\cong L^2(C_i,E\vert_{C_i})$ is compact.
But modulo the isomorphism $L^2(G,E_{i})^{G_{x_i}} \cong L^2(C_i,E\vert_{C_i})$ the operator $\hat{T}_\varphi^{G_{i}}$ is exactly $T_\varphi$,
and this shows that $B \circ T_\varphi$ is compact.
Since the ideal of compact operators is closed, we see that $M_{\phi_i} \maD M_{\phi_i}$ is compact.

Therefore, it follows that $\oid + \maD$ is Fredholm if, and only if, 
$\oid + p \maD p$ is Fredholm.
The operator $\oid + p \maD p$ is diagonal with respect to the orthogonal decomposition 
$L^2(M,E) =L^2(M\smallsetminus \maP , E\vert_{M\smallsetminus \maP}) \oplus L^2(\maP,E\vert_\maP)$.
Thus, $\oid + \maD$ is Fredholm if, and only if, 
$\oid_{L^2(M\smallsetminus \maP,E\vert_{M\smallsetminus \maP})} + p \maD p$ is Fredholm.

Since for any $\phi, \psi \in \maC(M\smallsetminus \maP)^G$ with disjoint support, 
\begin{equation*}
M_\phi\big(\oid_{L^2(M\smallsetminus \maP,E\vert_{M\smallsetminus \maP})} + p \maD p\big) M_\psi=M_\phi \maD M_\psi 
\end{equation*}
is a compact operator, we get that
\begin{equation*}
\oid_{L^2(M\smallsetminus \maP,E\vert_{M\smallsetminus \maP})} + p \maD p \in \Psi_{M\smallsetminus \maP/G}\big(L^2(M\smallsetminus \maP,E\vert_{M\smallsetminus \maP})\big).
\end{equation*}  
Consequently, the result follows from Proposition \ref{prop.gen.Simonenko}
because on $M\smallsetminus \maP$ the property of strong convergence to $0$ is satisfied. 
\end{proof}

\end{document}
